\documentclass[10pt,a4paper]{amsart}
\usepackage[margin=19mm]{geometry}
\usepackage{amsmath,amssymb,mathtools}
\usepackage[scr=rsfs,cal=euler]{mathalfa}
\usepackage{xfrac}
\usepackage[unicode,hidelinks,bookmarks=false]{hyperref}
\newcommand{\ie}{i.e.\ }
\newcommand{\cf}{cf.\ }
\newcommand{\resp}{respectively\ }
\newcommand{\Subsection}{\subsection}
\providecommand{\nobreakdash}{\nobreak\hbox{-}}
\setlength{\emergencystretch}{2em}
\multlinegap0pt
\usepackage{xfrac}
\usepackage{amssymb}
\usepackage{dsfont}
\usepackage{enumerate}
\usepackage[shortlabels]{enumitem}
\newtheorem{theorem}{Theorem}
\newtheorem{lemma}{Lemma}
\newtheorem{proposition}{Proposition}
\providecommand{\Dtil}{\widetilde D}
\DeclareMathOperator{\E}{\mathds{E}}
\DeclareMathOperator{\N}{N}
\providecommand{\R}{\mathbb R}
\providecommand{\Wtil}{\widetilde W}
\renewcommand{\ge}{\geqslant} 
\renewcommand{\le}{\leqslant} 
\title{Limiting Law of Local Fields for the Mean Field Ghatak-Sherrington Model}
\author{Yunhui Chen$^\star$, Keaton Fierro$^\dagger$}
\date{Source: arXiv:2608.10515v2 (CC BY 4.0)}
\begin{document}
\maketitle

\noindent\emph{Source and licence.} Limiting Law of Local Fields for the Mean Field Ghatak-Sherrington Model. Version arXiv:2608.10515v2 (CC BY 4.0), distributed by arXiv under the Creative Commons Attribution 4.0 International licence, http://creativecommons.org/licenses/by/4.0/, as recorded in the arXiv licence metadata for this exact version and shown on the abstract page https://arxiv.org/abs/2608.10515v2. Full licence notice: \texttt{licenses/arxiv-2608.10515-CC-BY-4.0.txt}.\par\smallskip

\noindent\textit{Editorial scope.} Counted unit: the article's lemma lem:onsager (source0.tex lines 962--967). The article's complete original proof of it is delivered with it (source0.tex lines 969--1028, one \texttt{proof} environment of 60 lines). No mathematics is written, repaired or re-derived: the statement and the proof are the article's own lines, in the article's own notation, and no retained line is substituted or re-flowed.  Retained uncounted as the source-local context the proof needs: lines 425--427; lines 450--456; lines 816--821; lines 839--846.  Nothing was omitted from any retained span, so the package carries no omission record.  No retained line contains a citation, so the wrapper carries no bibliography and no bibliography entry.  No retained line contains a figure environment, an image-inclusion command or drawing code, so no image is delivered and the package declares no asset dependency.  Editorial formatting only, in addition: an AMS \texttt{amsart} preamble that restates the article's own declarations and macros used by the retained lines, namely the environments \texttt{theorem}, \texttt{lemma}, \texttt{proposition} on separate counters, together with the standard packages those lines need.  The retained environments print in this document as Theorem 1, Lemma 1, Proposition 1 in order of appearance, whereas the article prints the same items with the numbering of its own full document.  The complete native source of the article as distributed in the arXiv e-print for this exact version is bundled unchanged as \texttt{sources/207331/source0.tex}.
\begin{equation}\label{eq:known-overlap-estimate}
  \E\langle |R_{1,2}-q|^{2k}\rangle \le \frac{C}{N^k}, \qquad \E\langle |R_{1,1}-p|^{2k}\rangle\le \frac{C}{N^k}.
\end{equation}

\begin{theorem}[CLT of cavity field]\label{thm:cavity}
Let $\beta_0$ be as above and $k\in\N$. Suppose that $U$ is an infinitely differentiable function on $\R$ whose derivatives of all orders are of moderate growth. Then for $\beta\le \beta_0$ and $D,h\in\R$, we have
\begin{equation}\label{eq:thm-cavity}
  \E\left(\langle U(\ell)\rangle-\int_{\R}U(x)\varphi_{r,p-q}(x)d x\right)^{2k}\le \frac{C}{N^k},
\end{equation}
where $C$ is a constant depending on $k,U,S,D,h$, and the high-temperature bound only. The second subscript in $\varphi_{u,v}$ is the variance, as in the convention above.
\end{theorem}

\begin{lemma}\label{lem:param-comparison}
There exists a positive constant $C$ independent of $N$ such that
\begin{equation}\label{eq:param-comparison}
  |\beta-\beta_-|,  \ |p-p_-|, \ |q-q_-|, \ |\Dtil-\Dtil_-|, \ |v-v_-|\le \frac{C}{N}.
\end{equation}
\end{lemma}

\begin{proposition}\label{prop:ratio}
Let $\beta\le \beta_0$ and $h,D\in\R$. Let $U:\R\to\R$ be smooth, with derivatives of moderate growth. For a cavity field $l$ and $r=\langle l\rangle$, with Gaussian reference variance $p-q$, we have, for every positive integer $k$,
\begin{equation}\label{eq:ratio-estimate}
  \E\left[\frac{\langle U(l)W(\beta l+h)\rangle}{\langle W(\beta l+h)\rangle}
  -\frac{\E_\xi[U(\xi+r)W(\beta(\xi+r)+h)]}{\Wtil(\beta r+h)}\right]^{2k}\le \frac{K}{N^k}.
\end{equation}
Here $K$ is uniform over the system size and over $0<\beta\le\beta_0$, by the uniformity of \eqref{eq:known-overlap-estimate}.
\end{proposition}

\begin{lemma}\label{lem:onsager}
For every $k\ge1$,
\begin{equation}\label{eq:onsager-bound}
  \E|r_- -\gamma_N|^{2k}\le \frac{C}{N^k}.
\end{equation}
\end{lemma}

\begin{proof}
Let
\[
  F_-(x):=\frac{\Wtil_-'(\beta_-x+h)}{\Wtil_-(\beta_-x+h)}.
\]
Since the spins are bounded by $S$, we have $|F_-(x)|\le S$. The cavity decomposition gives
\[
  \langle\sigma_N\rangle
  =\frac{\langle W'(\beta l+h)\rangle_-}{\langle W(\beta l+h)\rangle_-}
  =\frac{\langle W'(\beta_- l_-+h)\rangle_-}{\langle W(\beta_- l_-+h)\rangle_-},
\]
and
\[
  \langle l_N\rangle
  =\sqrt{\frac{N-1}{N}}\,
  \frac{\langle l_-W(\beta_-l_-+h)\rangle_-}{\langle W(\beta_-l_-+h)\rangle_-}.
\]
Apply Proposition~\ref{prop:ratio} to the $(N-1)$-spin system, first with
\[
  U_1(x):=\frac{W'(\beta_-x+h)}{W(\beta_-x+h)},
\]
for which $U_1(x)W(\beta_-x+h)=W'(\beta_-x+h)$ and
\[
  \frac{\E_{\xi_-}W'(\beta_-(r_-+\xi_-)+h)}{\Wtil_-(\beta_-r_-+h)}
  =F_-(r_-),
\]
and then with $U(x)=x$. Using Theorem~\ref{thm:cavity} inside the ratio estimate, there exist error variables $\Delta_1$ and $\Delta_2$ such that
\begin{equation}\label{eq:delta1}
  \langle\sigma_N\rangle=F_-(r_-)+\Delta_1,
  \qquad
  \E|\Delta_1|^{2k}\le \frac{C}{N^k},
\end{equation}
and
\begin{equation}\label{eq:delta2}
  \frac{\langle l_-W(\beta_-l_-+h)\rangle_-}{\langle W(\beta_-l_-+h)\rangle_-}
  =r_-+\beta_-v_-F_-(r_-)+\Delta_2,
  \qquad
  \E|\Delta_2|^{2k}\le \frac{C}{N^k}.
\end{equation}
The second identity uses the Gaussian integration-by-parts formula
\[
  \E_{\xi_-}[(r_-+\xi_-)W(\beta_-(r_-+\xi_-)+h)]
  =r_-\Wtil_-(\beta_-r_-+h)+\beta_-v_-\Wtil_-'(\beta_-r_-+h).
\]
Consequently, writing $c_N=\sqrt{(N-1)/N}$, equations \eqref{eq:delta1} and \eqref{eq:delta2} give
\begin{equation}\label{eq:gamma-r-explicit}
  \gamma_N-r_-=(c_N-1)r_-+(c_N\beta_-v_- -\beta v)F_-(r_-)+c_N\Delta_2-\beta v\Delta_1.
\end{equation}
By Lemma~\ref{lem:param-comparison},
\[
  |c_N-1|+|c_N\beta_-v_- -\beta v|\le \frac{C}{N}.
\]
Moreover $r_-$ has Gaussian moments of all orders uniformly in $N$, since it is a Gaussian linear combination of bounded magnetizations. Since $|F_-(r_-)|\le S$, \eqref{eq:gamma-r-explicit}, the elementary inequality $|x_1+\cdots+x_4|^{2k}\le C_k\sum_i|x_i|^{2k}$, and \eqref{eq:delta1}-\eqref{eq:delta2} imply
\[
  \E|\gamma_N-r_-|^{2k}
  \le C\left(\frac{\E|r_-|^{2k}}{N^{2k}}+\frac1{N^{2k}}+\E|\Delta_1|^{2k}+\E|\Delta_2|^{2k}\right)
  \le \frac{C}{N^k}.
\]
This proves the lemma.
\end{proof}

\end{document}
